%%% FIELD fixed-graph-reduction-proof
Assume
\[
\mathfrak G_{H,L}(\mathcal S)=\{G_0\}
\qquad\text{and}\qquad
\mathfrak F(\mathcal S,L,H,p)\ne\varnothing.
\]
Define the canonical feasible set
\[
K_0(p)=\{\lambda\in\Lambda_{G_0}:P_{G_0}(\lambda)=p\}.
\]
The second premise gives a pair \((G_0,\mathcal M)\in\mathfrak F(\mathcal S,L,H,p)\), because the first premise makes \(G_0\) the only admissible refinement.  By the compatible-SCM-to-canonical-representative assertion of \(\text{thm:sharp-refinement-envelope}\), there is a \(\lambda\in\Lambda_{G_0}\) satisfying \(P_{G_0}(\lambda)=p\) and \(q_{G_0}(\lambda)=q(\mathcal M)\).  Hence
\[
K_0(p)\ne\varnothing.
\]
It follows that the set of refinements retained by the feasible-refinement restriction in \(\text{thm:sharp-refinement-envelope}\) is
\[
\bigl\{G\in\mathfrak G_{H,L}(\mathcal S):
\{\lambda\in\Lambda_G:P_G(\lambda)=p\}\ne\varnothing\bigr\}
=\{G_0\}.
\]
Substitution of this equality into the two outer extrema in that theorem yields, term by term,
\[
\underline q(p)
=\min_{\lambda\in K_0(p)}q_{G_0}(\lambda)
=\min_{\lambda\in\Lambda_{G_0}:P_{G_0}(\lambda)=p}q_{G_0}(\lambda)
\]
and
\[
\overline q(p)
=\max_{\lambda\in K_0(p)}q_{G_0}(\lambda)
=\max_{\lambda\in\Lambda_{G_0}:P_{G_0}(\lambda)=p}q_{G_0}(\lambda).
\]
These are exactly the objective and equality constraints of the single \(G_0\)-program in \(\text{def:raw-polynomial-compiler}\).  The minima and maxima exist, rather than merely being infima and suprema, by the attainment conclusion of \(\text{thm:sharp-refinement-envelope}\).

Now also assume \(E_{G_0}^{\leftrightarrow}=\varnothing\).  Because \(G_0\) is an ADMG, deleting all bidirected edges leaves its acyclic directed part, so \(G_0\) is a DAG.  By the definition of \(\Lambda_{G_0}\), the exogenous-parent indices are then only the singletons \(a=\{v\}\), and therefore
\[
\Lambda_{G_0}
=\bigsqcup_{r\in\mathcal R_{G_0}}
\{r\}\times
\prod_{v\in\mathcal W_H}
\Delta(\mathcal K_{G_0,r,\{v\}}).
\]
Fix a branch \(r\), abbreviate \(\lambda_v=\lambda_{\{v\}}\), and let
\(\kappa=(\kappa_v)_{v\in\mathcal W_H}\) be a tuple of private response states.  Mutual independence of the exogenous parents, retained in the product-of-simplices parameterization used by \(\text{thm:sharp-refinement-envelope}\), gives the tuple mass
\[
\mu_{r,\lambda}(\kappa)
=\prod_{v\in\mathcal W_H}\lambda_v(\kappa_v).
\]
By \(\text{ass:structural-recursion}\), the fixed table branch recursively determines an observational assignment \(w^r(\kappa)\) and a deterministic interventional contrast \(d^r(\kappa)\in\{-1,0,1\}\).  Consequently the compiled coordinates are
\[
P_{G_0}(\lambda)(w)
=\sum_{\kappa}
\mathbf 1\{w^r(\kappa)=w\}
\prod_{v\in\mathcal W_H}\lambda_v(\kappa_v)
\]
for every \(w\in\{0,1\}^{\mathcal W_H}\), and
\[
q_{G_0}(\lambda)
=\sum_{\kappa}d^r(\kappa)
\prod_{v\in\mathcal W_H}\lambda_v(\kappa_v).
\]
Thus every observational constraint and the causal objective is a sum of products of independent node-specific response weights, and each structural table uses only the observed DAG parents of \(v\) and the private response state \(\kappa_v\).  This is precisely the independent node-response-factor program for the known DAG.  Zero simplex weights are allowed, so this reduction uses neither strict positivity nor division.
